\section{Service Architecture Used}
\label{sec:architecture}

To satisfy the scalability, observability, and threat-resilience mandates of Project 1, I implemented a modular, 7-tier Service-Oriented Architecture deployed across isolated Linux Containers (Incus). The complete architectural topology and data flow are illustrated in \cref{fig:architecture}.

\begin{figure}[H]
\centering
\begin{tikzpicture}[
    scale=0.80, transform shape,
    service/.style={rectangle, draw=blue!70!black, fill=blue!8, rounded corners, thick, minimum width=2.6cm, minimum height=1.1cm, align=center, font=\footnotesize\bfseries},
    edge_node/.style={rectangle, draw=purple!70!black, fill=purple!8, rounded corners, thick, minimum width=2.4cm, minimum height=1.1cm, align=center, font=\footnotesize\bfseries},
    storage/.style={cylinder, shape border rotate=90, draw=orange!80!black, fill=orange!12, aspect=0.25, thick, minimum width=2.4cm, minimum height=1.2cm, align=center, font=\footnotesize\bfseries},
    security/.style={rectangle, draw=red!70!black, fill=red!8, rounded corners, thick, minimum width=3.0cm, minimum height=1.1cm, align=center, font=\footnotesize\bfseries},
    monitor/.style={rectangle, draw=green!60!black, fill=green!8, rounded corners, thick, minimum width=2.8cm, minimum height=1.1cm, align=center, font=\footnotesize\bfseries},
    edge/.style={-Latex, thick, draw=black!75},
    secedge/.style={-Latex, dashed, thick, draw=red!70!black},
    regedge/.style={-Latex, dashed, thick, draw=blue!70!black},
    monedge/.style={-Latex, thick, draw=green!60!black}
]

% ==================== ROW 1 (TOP): Service Discovery ====================
\node[service, minimum width=4.0cm] (registry) at (5.7, 2.5) {Service Registry (\texttt{:5050})\\{\scriptsize\normalfont Token Auth $\cdot$ Rate Limit $\cdot$ Agent}};

% ==================== ROW 2 (MIDDLE): Data Pipeline ====================
\node[edge_node] (sensor) at (0.0, 0.0) {Edge Sensor (\texttt{:5000})\\{\scriptsize\normalfont IoT Telemetry $\cdot$ HMAC Sign}};

\node[service] (ingest) at (3.8, 0.0) {Ingestion Gateway (\texttt{:5001})\\{\scriptsize\normalfont DoS Shield $\cdot$ Access Log $\cdot$ Agent}};

\node[storage] (storage) at (7.6, 0.0) {Storage Service (\texttt{:5002})\\{\scriptsize\normalfont SQLite WAL $\cdot$ Access Log $\cdot$ Agent}};

\node[service] (analytics) at (11.4, 0.0) {Analytics Worker (\texttt{:5003})\\{\scriptsize\normalfont Math Stats $\cdot$ Access Log $\cdot$ Agent}};

% ==================== CLUSTER BOUNDARY (Incus Containers) ====================
\begin{scope}[on background layer]
\node[draw=blue!50!black, dashed, fill=blue!2, rounded corners=8pt,
      fit=(registry) (ingest) (storage) (analytics),
      inner sep=10pt,
      label={[anchor=north west, font=\scriptsize\bfseries\color{blue!80!black}, xshift=6pt, yshift=-3pt]above left:Incus Zero-Trust Microservice Cluster (wazuh-agent + iptables on each node)}] (cluster) {};
\end{scope}

% ==================== ROW 3 (BOTTOM): Exporters & SIEM ====================
\node[security] (wazuh) at (1.6, -3.0) {Wazuh SIEM Manager (\texttt{:1514})\\{\scriptsize\normalfont Rules Engine $\cdot$ Active Response}};

\node[monitor] (diag) at (5.7, -3.0) {Diagnostics Exporter (\texttt{:8001})\\{\scriptsize\normalfont RPS, CPU, Error Metrics}};

\node[monitor] (health) at (9.8, -3.0) {Health Exporter (\texttt{:8000})\\{\scriptsize\normalfont Periodic Liveness Poller}};

% ==================== ROW 4: Prometheus ====================
\node[monitor, minimum width=4.0cm] (prom) at (7.7, -4.9) {Prometheus Server (\texttt{:9090})\\{\scriptsize\normalfont TSDB Time-Series $\cdot$ PromQL Alerts}};

% ----------------- Connections: Data Flow -----------------
\draw[edge] (sensor) -- node[above=1pt, font=\tiny] {HMAC Sign} node[below=1pt, font=\tiny] {\texttt{POST /ingest}} (ingest);
\draw[edge] (ingest) -- node[above=1pt, font=\tiny] {\texttt{POST /store}} (storage);
\draw[edge] (storage) -- node[above=1pt, font=\tiny] {Read SQLite} (analytics);

% ----------------- Connections: Registry Discovery -----------------
\draw[regedge] (sensor.north) |- node[near start, left=2pt, font=\tiny] {Heartbeat} (registry.west);
\draw[regedge] (ingest.north) -- node[left=2pt, font=\tiny] {Resolve Storage} (registry.215);
\draw[regedge] (storage.north) -- node[right=2pt, font=\tiny] {Register} (registry.325);
\draw[regedge] (analytics.north) |- node[near start, right=2pt, font=\tiny] {Register} (registry.east);

% ----------------- Connections: Wazuh SIEM -----------------
\draw[secedge] (cluster.205) to[bend right=12] node[above left, font=\tiny, align=right] {All Access Logs\\(\texttt{/var/log/ssle/*})} (wazuh.60);
\draw[secedge] (wazuh.25) to[bend right=12] node[below right, font=\tiny, align=left] {Active Response: Cluster \texttt{firewall-drop}\\(\texttt{<location>all</location>})} (cluster.225);

% ----------------- Connections: Exporters -----------------
\draw[monedge] (diag.north west) -- node[above left=1pt, font=\tiny] {Runtime Diagnostics} (ingest.south east);
\draw[monedge] (health.north west) -- node[above left=1pt, font=\tiny] {Poll /health} (storage.south east);
\draw[monedge] (health.north east) -- node[above right=1pt, font=\tiny] {Poll /health} (analytics.south);

% ----------------- Connections: Prometheus Scraping -----------------
\draw[monedge] (prom.north west) -- node[below left=1pt, font=\tiny] {Scrape :8001} (diag.south);
\draw[monedge] (prom.north east) -- node[below right=1pt, font=\tiny] {Scrape :8000} (health.south);

\end{tikzpicture}
\caption{End-to-End SSLE Project 1 Architecture: Multi-tier microservices, dual Prometheus observability exporters, cluster-wide Wazuh SIEM log shipping, and automated Active Response feedback loop.}
\label{fig:architecture}
\end{figure}

\subsection{Microservice Functional Decoupling}
The architecture decomposes application concerns into distinct microservices:
\begin{enumerate}
    \item \textbf{Service Registry (\texttt{:5050}):} Acts as the authoritative, in-memory catalog of all running services. Maintains location metadata, enforces token-based administrative protection for service deregistration, supports location-filtered queries, and applies token-bucket rate limiting ($30\,\text{req/s}$) and payload clamping to prevent control-plane starvation.
    \item \textbf{Edge Temperature Sensor (\texttt{:5000}):} Simulates IoT edge sensor hardware. Emits periodic temperature telemetry samples, dynamically computes HMAC-SHA256 integrity signatures, attaches Unix epoch timestamps, and queries the registry to locate the ingestion gateway.
    \item \textbf{Hardened Ingestion Gateway (\texttt{:5001}):} The perimeter microservice facing external networks. Enforces defense-in-depth: token-bucket rate limiting ($5\,\text{req/s}$ sustained, $15\,\text{req/s}$ burst threshold), HTTP payload size clamping ($1\,\text{MB}$ via \texttt{http.MaxBytesReader} returning \texttt{HTTP 413}), Shellshock CVE-2014-6271 header inspection, and cryptographic HMAC signature and timestamp replay validation. Outputs structured Apache Combined access logs to \texttt{/var/log/ssle/}\allowbreak\texttt{ingestion\_access.log} for SIEM processing.
    \item \textbf{Storage Service (\texttt{:5002}):} Interacts with an embedded SQLite relational database located on a mounted storage volume (\texttt{/data/temp\_data.db}), providing indexed schema persistence for time-series temperature tuples, protected by local rate limiting ($30\,\text{req/s}$) and structured access logging.
    \item \textbf{Analytics Service (\texttt{:5003}):} Reads persisted readings from the database, computes sliding-window statistical metrics (mean temperature), exposes analytical endpoints (\texttt{/average}), and enforces access logging and rate limiting.
    \item \textbf{Prometheus Health Exporter (\texttt{:8000}):} Periodically polls liveness and readiness endpoints of all microservices, attaching dimensional location labels to exportable Prometheus gauges.
    \item \textbf{Prometheus Diagnostics Exporter (\texttt{:8001}):} Instruments runtime operational performance: requests per second (RPS), CPU utilization percent, memory allocations, and HTTP error rate percentages.
\end{enumerate}

\subsection{Cluster-Wide Defense-in-Depth and Zero-Trust Hardening}
To reflect realistic production environments where an internal attacker or compromised workload could pivot laterally across private container networks, security controls are not restricted solely to the public-facing ingestion gateway. I architected defense-in-depth across both the application middleware and host network layers:
\begin{itemize}
    \item \textbf{Universal Application-Level DoS Middleware:} All core Go microservices (Ingestion, Registry, Storage, and Analytics) integrate the shared \texttt{DoSProtectionMiddleware}. It provides local token-bucket rate limiting and enforces strict $1\,\text{MB}$ body size clamping returning \texttt{HTTP 413 Payload Too Large}. This prevents rogue or malfunctioning internal services from overwhelming downstream databases or exhausting connection pools.
    \item \textbf{Decentralized Access Logging for SIEM Ingestion:} Every core microservice produces structured Apache Combined access logs written to dedicated paths in \texttt{/var/log/ssle/}\allowbreak\texttt{<service>\_access.log}. This ensures that every HTTP transaction, status code, and client IP across the entire cluster is auditable.
    \item \textbf{Host-Level Packet Filtering (\texttt{iptables}):} Each microservice container runs isolated kernel firewall rules configured during deployment. Inbound policies drop unsolicited traffic, permitting only loopback, established connections, and the designated TCP listening port for that service.
    \item \textbf{Distributed Wazuh Agents and Active Response Feedback:} Each core container runs an active \texttt{wazuh-agent} daemon that forwards local access logs to the centralized Wazuh Manager. When an attack signature or repeated threshold breach is detected, the manager broadcasts active response \texttt{firewall-drop} commands across all microservice containers simultaneously, isolating the attacker IP globally.
\end{itemize}

\subsection{Hierarchical Identification and Fully Qualified Service Names (FQSN)}
A core requirement of Project 1 is the identification and registration of elements in a scalable, location-aware structure that encodes physical hardware, logical container virtualization, and geographic delegation of responsibility.

To eliminate namespace collisions and enable instant incident localization, I formulated the \textbf{Fully Qualified Service Name (FQSN)} schema:
\begin{equation}
    \text{FQSN} = \langle\text{region}\rangle.\langle\text{zone}\rangle.\langle\text{physical\_host}\rangle.\langle\text{logical\_env}\rangle.\langle\text{instance\_id}\rangle.\langle\text{kind}\rangle
\end{equation}
A concrete instantiation of this naming standard in the testbed is:
\begin{lstlisting}[language=bash]
pt-aveiro.lab-ssle.kali-vm.incus:ssle-ingestion.a1b2c3d4.ingestion
\end{lstlisting}
The semantic components of this hierarchy provide distinct operational capabilities:
\begin{itemize}
    \item \textbf{Geographic Administrative Boundary (\texttt{pt-aveiro.lab-ssle}):} Designates regional authority, regulatory jurisdiction (e.g., GDPR data residency), and network routing boundaries.
    \item \textbf{Physical Infrastructure Layer (\texttt{kali-vm}):} Identifies the bare-metal hypervisor, physical blade server, or host machine.
    \item \textbf{Logical Virtualization Layer (\texttt{incus:ssle-ingestion}):} Identifies the isolated Linux container, profile, and cgroup boundary.
    \item \textbf{Ephemeral Instance Identifier (\texttt{a1b2c3d4}):} A cryptographically random hex UUID ensuring idempotency and collision avoidance during autoscaling.
    \item \textbf{Functional Service Class (\texttt{ingestion}):} Defines the functional interface contract and routing target.
\end{itemize}

\subsection{Rapid Incident Localization API}
When an alert fires in production, resolving an alert to the exact container and physical host is critical for Mean Time to Repair (MTTR). The Service Registry implements a specialized localization endpoint:
\begin{equation}
    \texttt{GET /services/locate/\{query\}}
\end{equation}
where \texttt{\{query\}} accepts an instance ID, exact/partial FQSN, endpoint URL, or functional kind. The registry resolves the query in $O(1)$ time, returning a structured remediation payload (\cref{lst:locate}).

\begin{lstlisting}[caption={Rapid Incident Localization Response for \texttt{ingestion}.}, label={lst:locate}]
{
  "status": "FOUND",
  "query": "ingestion",
  "incident": {
    "fqsn": "pt-aveiro.lab-ssle.kali-vm.incus:ssle-ingestion.9494d3aa.ingestion",
    "physical_host": "kali-vm",
    "logical_env": "incus:ssle-ingestion",
    "administrative_domain": "pt-aveiro.lab-ssle",
    "url": "http://ssle-ingestion.incus:5001",
    "remediation_hint": "Run to inspect container: incus exec ssle-ingestion -- journalctl -u ssle-ingestion -n 50"
  }
}
\end{lstlisting}
This allows monitoring engines, operators, and automated SIEM response scripts to instantly identify the host to inspect or container to isolate without manual cross-referencing.
